\scanpage{209}

\uline{Question} Soit $E$ module à connexion \deleted{intégrable}\inserted{à courbure nulle} sur $X/S$ en car. $p>0$. Peut-on définir un tenseur \deleted{\unsure{dit de}} $p$-courbure sur $X$, qui soit

\[
K\in\Gamma(X,\underline{\mathrm{End}}(E)\otimes(\Omega^{\deleted{\unsure{\otimes p}}1}_{X/S})^{\otimes p})\qquad ?
\]

On a aussi une \deleted{\unsure{hom|homomorphie}} suite exacte, où $\underline{\mathrm{Pic}}^{ci}_{X/S}$ désigne l'image de \deleted{$\unsure{\mathrm{Pic}}$} $P^{ci}_{X/S}\longrightarrow\underline{\mathrm{Pic}}_{X/S}$ (cf. (2.2) dans le cas lisse)

\[
0\longrightarrow
\underbrace{[f_*^{\mathrm{fl}}(\underline O^*_{X/S})\longrightarrow \overset{\inserted{\substack{\text{(défaut d'exactitude}\\\text{dans le milieu)}}}}{f_*^{\mathrm{fl}}(\underline Z^1_{X/S})}\longrightarrow f_*^{\mathrm{fl}}(\underline\Omega^{1\otimes p}_{X/S})]}_{\underline{\mathrm{Cart}}_{X/S}\ \text{``faisceau de Cartier''}}
\longrightarrow P^{cip}_{X/S}\longrightarrow\underline{\mathrm{Pic}}^{ci}_{X/S}
\longrightarrow f_*^{\mathrm{fl}}(\Omega^{1\otimes p}_{X/S})/\mathrm{Im}\,f_*^{\mathrm{fl}}(\underline Z^1_{X/S}),
\tag*{(3.2)}
\]

où la flèche

\[
K_0\colon f_*^{\mathrm{fl}}(\underline Z^1_{X/S})\longrightarrow f_*^{\mathrm{fl}}(\underline\Omega^{1\otimes p}_{X/S})
\]
est donnée par

\[
K_0(\omega)=\omega^{\otimes p}-{\unsure{F|f}}_{X/S}^*(\underbrace{\widetilde C}_{\text{\inserted{opération de Cartier}}}\omega)\qquad\text{pour $d\omega=0$,}\tag*{(3.3)}
\]

\uline{NB} $K_0$ \inserted{est additif mais} n'est pas linéaire, c'est la somme d'une opération $p$-linéaire et d'une opération linéaire.

Tjrs dans le cas de l'existence locale d'une $p$-base, on sait qu'\deleted{\unsure{un}}\inserted{pour un} module \deleted{inversible} qu. coh. $E$ sur $X$, la donnée d'une connexion à courbure et $p$-courbure nulles \deleted{\unsure{éq}} sur $E$ équivaut à celle d'une \uline{donnée} \uline{de descente} de $E$ à $\Lambda=\deleted{\underline O_{\unsure{S|X}}^p}\underline O_S[\underline O_X^p]\simeq\underline O_{X^{p/S}}/J$, où $J$ est le noyau de $\underline O_{X^{p/S}}\longrightarrow F_{X/S\,*}(\underline O_X)$,

\deleted{et qu'une telle donnée de des.} \inserted{La catégorie de ces} $E$ équivaut à celle des modules quasi-cohérents sur $(X^{p/S},\Lambda)\simeq(X,\underline O_X^p)$. Soit $(X^{p/S\,!})$ cet espace annelé, on trouve donc que $P^{cip}_{X/S}$ est le faisceau associé au préfaisceau

\[
S'\longmapsto\mathrm{Pic}(X_{S'}^{p/S'\,!}).
\]

% archive sheet number: 208
